\section{从模仿到优化：RLHF 的动机}

\subsection{范式转变：从生成模型到策略优化}

\begin{figure}[H]
\centering
\includegraphics[width=0.95\textwidth]{slides-images/slide-030.jpg}
\caption{两种范式的对比——模仿学习（SFT）vs. 优化（RLHF）\protect\footnotemark}
\end{figure}
\footnotetext{Slides 第31页。}

SFT 和 RLHF 代表了两种根本不同的视角：

\[
\text{SFT（模仿）: } \hat{p}(y|x) \approx p^*(y|x)
\]

其中 $p^*$ 是某个参考分布（专家示范数据的分布）。这是纯粹的生成建模视角——我们试图让模型分布逼近参考分布。

\[
\text{RLHF（优化）: }\quad
\text{找到 } \hat{p}(y|x) \text{ 使得 }
\max_p \mathbb{E}_p[R(y,x)]
\]

其中 $R(y,x)$ 是奖励函数。在这个视角下，语言模型不再是某个分布的模型，而是需要最大化奖励的\textbf{策略}。

\begin{importantbox}{为什么要做 RLHF？两个核心原因}
\begin{enumerate}
    \item \textbf{数据效率}：SFT 需要专家编写完整的回复（昂贵且费时），而 RLHF 只需要对现有回复做成对比较（更快、更便宜）
    \item \textbf{生成-验证差距}：人类判断回复好坏的能力，往往优于自己生成高质量回复的能力
\end{enumerate}
\end{importantbox}

\subsection{SFT 数据收集的高成本}

\begin{figure}[H]
\centering
\includegraphics[width=0.95\textwidth]{slides-images/slide-032.jpg}
\caption{后训练流程中不同阶段的成本分布：SFT 的人工数据收集成本远高于 RLHF 的成对比较数据\protect\footnotemark}
\end{figure}
\footnotetext{Slides 第33页。}

前沿实验室在后训练数据上的花费可达数百万美元。SFT 数据需要专家撰写长文本回复——正如课堂实验所展示的，即便是高度积极的斯坦福学生，在短时间内也很难写出高质量的长回复。相比之下，``A 比 B 好''这样的成对判断要容易得多。

\subsection{生成-验证差距}

一项来自 Tatsu 实验室的研究发现了一个令人惊讶的现象：自由撰稿专家在被要求写摘要后，评估时竟然更偏好 AI 生成的摘要而非自己写的。原因在于人类写作时会受到各种约束（时间压力、``应该用华丽语言''的心理预设等），但评估时却能更客观地判断质量。

\begin{knowledgebox}{验证比生成更容易也更准确}
这一发现（验证比生成更容易）不仅是 RLHF 的动机之一，也是计算复杂性理论中的经典命题（P vs NP）在实际场景中的体现。在后训练语境下，它意味着：利用人类的验证能力（通过成对比较）可能比利用人类的生成能力（通过撰写回复）更加高效和准确。
\end{knowledgebox}


